\section{Kapitel 4 --
Systemkonzeption}\label{kapitel-4-systemkonzeption}

\begin{center}\rule{0.5\linewidth}{0.5pt}\end{center}

\subsection{4.1 Überblick und
Schichtenmodell}\label{uxfcberblick-und-schichtenmodell}

Das System ist als \textbf{eigenständiger Cloud-Service} konzipiert, der
alle Spiele der konfigurierten Wettbewerbe automatisch verarbeitet ---
von der Datenimportierung über die KI-Textgenerierung bis zur
Publikation. Es ist eine serviceorientierte, dreischichtige Architektur,
die Datenbeschaffung, Anwendungslogik und Präsentation klar voneinander
trennt. Diese Schichtentrennung folgt dem etablierten Prinzip der
\emph{Separation of Concerns} (Parnas 1972), das die unabhängige
Weiterentwicklung einzelner Systemteile ermöglicht.

Die \textbf{Datenbeschaffungs- und Automatisierungsschicht} wird über
n8n umgesetzt. Sie integriert externe Datenquellen (insbesondere
Football-API sowie Medienquellen wie ScorePlay und Social-Feeds),
transformiert die Rohdaten in ein einheitliches internes Format und
schreibt diese strukturiert in die Persistenzschicht beziehungsweise
über definierte Backend-Endpunkte in das System ein. Die Entkopplung der
Integrationslogik in n8n ermöglicht eine schnelle Anpassung von
Import-Workflows, ohne den Anwendungskern ändern zu müssen.

Die \textbf{Anwendungs- und Persistenzschicht} basiert auf FastAPI
(Python) und PostgreSQL. Sie stellt die REST-Schnittstellen bereit,
verwaltet den Zustandsübergang der Ticker-Einträge (z. B. draft,
published, rejected), orchestriert KI-Generierungsschritte und übernimmt
die konsistente Speicherung aller Domänendaten (u. a. Teams,
Wettbewerbe, Spiele, Events, Statistiken, Ticker- und Media-Daten).

Die \textbf{Präsentationsschicht} ist als White-Label-Frontend in React
und TypeScript realisiert. Sie bildet den Redaktionsworkflow für
Auswahl, Sichtung, Bearbeitung und Publikation ab und bleibt dabei
mandantenfähig konfigurierbar (z. B. Instanz, Stil, Sprache, Branding).
Technisch kombiniert das Frontend REST-basiertes Polling für Ticker- und
Spieldaten mit einem dedizierten WebSocket-Kanal für
Echtzeit-Medienupdates.

Die folgende Abbildung zeigt die drei Schichten mit ihren Technologien
und Kommunikationspfaden:

\begin{lstlisting}
graph LR
    classDef schicht fill:#EFF6FF,stroke:#3B82F6,stroke-width:1.5px,color:#1E3A8A
    classDef extern fill:#F9FAFB,stroke:#9CA3AF,stroke-width:1px,color:#374151
    classDef db fill:#F0FDF4,stroke:#22C55E,stroke-width:1.5px,color:#14532D

    subgraph EXT['  Externe Dienste  ']
        FAPI['Football-API\napi-sports.io']
        LLM['LLM Provider\nOpenRouter']
        SP['ScorePlay\nMedia']
        EF['EF Spielerdaten\nprofis.eintracht.de']
    end

    subgraph AUTO['  Automatisierungsschicht  ']
        N8N['n8n\n15 Workflows']
    end

    subgraph APP['  Anwendungs- & Persistenzschicht  ']
        BE['FastAPI\nPython · Uvicorn']
        DB[('PostgreSQL\n18 Tabellen')]
        BE <-->|'ORM'| DB
    end

    subgraph UI['  Präsentationsschicht  ']
        FE['React Frontend\nTypeScript']
    end

    N8N -->|'HTTP GET'| FAPI
    N8N -->|'HTTP GET/POST'| SP
    N8N -->|'HTTP GET'| EF
    N8N -->|'Zusammenfassungen'| LLM
    N8N -->|'REST POST /api/v1'| BE

    BE -->|'LLM-Aufrufe'| LLM

    FE -->|'REST Polling · 5s'| BE
    FE -.->|'WebSocket /ws/media'| BE
    FE -->|'Webhook-Trigger'| N8N

    class FE schicht
    class BE schicht
    class N8N schicht
    class DB db
    class FAPI,LLM,SP,EF extern
\end{lstlisting}

\begin{center}\rule{0.5\linewidth}{0.5pt}\end{center}

\subsubsection{4.1.1 Kommunikations- und
Triggerarchitektur}\label{kommunikations--und-triggerarchitektur}

Der Kommunikationsfluss folgt dem Prinzip \textbf{``read first, trigger
if missing''}: Das Frontend liest zunächst den vorhandenen Datenbestand
über REST; nur bei fehlenden Daten wird ein passender n8n-Webhook
ausgelöst. Dadurch werden externe API-Aufrufe minimiert und redundante
Importe vermieden.

\begin{lstlisting}
sequenceDiagram
    participant FE as React Frontend
    participant BE as FastAPI Backend
    participant N8N as n8n
    participant FAPI as Football-API

    FE->>BE: GET /teams/countries
    alt Keine Daten vorhanden
        FE->>N8N: Webhook import-countries
        N8N->>FAPI: GET /countries
        N8N->>BE: POST /api/v1/countries (Upsert)
        FE->>BE: GET /teams/countries (erneut)
    end

    FE->>BE: GET /matches (nach Teamauswahl)
    FE->>N8N: Webhook import-prematch, import-lineups, ...
    N8N->>BE: POST /api/v1/events, /api/v1/ticker/generate/...

    loop Polling alle 5s (Events, Ticker, Match)
        FE->>BE: GET /api/v1/ticker/{match_id}
        BE-->>FE: TickerEntry[] (inkl. neuer Entwürfe)
    end

    N8N->>BE: POST /api/v1/media/incoming
    BE-->>FE: WebSocket /ws/media → neue Medien
\end{lstlisting}

Die Systemkopplung erfolgt über drei klar getrennte Schnittstellentypen:

\begin{itemize}

\item
  \textbf{REST-API} --- stabiler Kern für alle Lese- und
  Schreiboperationen: Navigation, Ticker-CRUD, Statusübergänge
  (\passthrough{\lstinline!draft!} → \passthrough{\lstinline!published!}
  → \passthrough{\lstinline!rejected!}), manuelle Eingaben.
\item
  \textbf{n8n-Webhooks} --- bedarfsgesteuerte Trigger für drei Klassen:

  \begin{enumerate}
  \def\labelenumi{\arabic{enumi}.}
  
  \item
    \textbf{Initialisierungs-Trigger} --- Aufbau fehlender Stammdaten
    (Länder, Teams, Wettbewerbe, Spiele)
  \item
    \textbf{Match-Kontext-Trigger} --- spielbezogene Importe nach
    Matchauswahl (Events, Lineups, Statistiken, Prematch)
  \item
    \textbf{Generierungs-Trigger} --- Anstoß von KI-Textprozessen
    (Event-basiert, Matchphasen, Zusammenfassungen)
  \end{enumerate}
\item
  \textbf{WebSocket \passthrough{\lstinline!/ws/media!}} ---
  Echtzeit-Push für latenzkritische Medieninhalte (ScorePlay-Bilder,
  Social-Media-Clips); Ticker- und Matchdaten bleiben bewusst im
  Polling-Modell.
\end{itemize}

Diese hybride Triggerarchitektur unterstützt sowohl redaktionelle
Kontrolle im Sinne eines \textbf{Human-in-the-Loop}-Ansatzes als auch
weitgehend automatisierte Live-Prozesse.

\begin{center}\rule{0.5\linewidth}{0.5pt}\end{center}

\subsubsection{4.1.2 Partner-Team-Konzept und
White-Label-Steuerung}\label{partner-team-konzept-und-white-label-steuerung}

Das System unterscheidet zwischen dem \textbf{Partner-Team} --- dem
Verein, für den die White-Label-Instanz konfiguriert ist --- und den
jeweiligen Gegnern. Diese Unterscheidung wird über ein konfigurierbares
Team-Keyword umgesetzt, das sowohl in den n8n-Workflows als auch im
Frontend gegen den Teamnamen abgeglichen wird.

Auf Basis dieser Erkennung werden zwei Verarbeitungspfade gesteuert:

\begin{itemize}

\item
  \textbf{Instanz und Stil}: Partner-Team-Spiele verwenden den
  \passthrough{\lstinline!ef\_whitelabel!}-Instanzpfad mit euphorischem
  Stil; alle anderen Spiele den \passthrough{\lstinline!generic!}-Pfad
  mit neutralem Stil.
\item
  \textbf{UI-Anpassung}: Das Frontend blendet vereinsspezifische
  Bereiche (Social-Media-Panels, Torjubel-Videos) ausschließlich bei
  erkannten Partner-Team-Spielen ein.
\item
  \textbf{Erweiterte Kontextdaten}: Für Tormomente des Partner-Teams
  werden zusätzlich Spielerdaten und Torjubel-Video-URLs aus einer
  externen Vereinsquelle in den Prompt-Kontext eingebettet.
\end{itemize}

Das Datenbankschema enthält ergänzend ein
\passthrough{\lstinline!is\_partner\_team!}-Flag, das für zukünftige
Erweiterungen vorgesehen ist. Die konkreten Implementierungen sind in
Kapitel 5.2.2 (Datenbankmodell,
\passthrough{\lstinline!is\_partner\_team!}-Flag) und Kapitel 5.4.6
(Frontend-Moduslogik, \passthrough{\lstinline!isOurTeam!}-Erkennung)
dokumentiert.

\begin{center}\rule{0.5\linewidth}{0.5pt}\end{center}

\subsection{4.2 Backend-Konzeption}\label{backend-konzeption}

\subsubsection{4.2.1 Framework-Wahl:
FastAPI}\label{framework-wahl-fastapi}

Als Backend-Framework wurde FastAPI gewählt (vgl. Abschnitt 3.6.2 für
die technische Einordnung). Die Entscheidung begründet sich durch drei
systemspezifische Anforderungen: Erstens erfordert die parallele
Verarbeitung von LLM-Aufrufen, externen API-Abfragen und
Medienverarbeitung native Unterstützung für asynchrone I/O. Zweitens
vereinfacht die automatisch generierte OpenAPI-Spezifikation die
Integration mit n8n-Workflows erheblich, da Endpunkte über Swagger UI
direkt testbar sind. Drittens reduziert die enge Pydantic-Integration
Schnittstellenfehler zwischen Frontend, Backend und n8n, indem alle
Eingaben gegen typisierte Schemas validiert werden.

\begin{center}\rule{0.5\linewidth}{0.5pt}\end{center}

\subsubsection{4.2.2 Interne Struktur und
Datenzugriff}\label{interne-struktur-und-datenzugriff}

Die interne Struktur folgt einer klaren Trennung aus API-Routern,
Repository-Schicht (vgl. Kap. 3.6.3) sowie \textbf{SQLAlchemy
2.0}-ORM-Modellen und Pydantic-Schemas. Die Router bilden die
HTTP-Schnittstelle und übernehmen Validierung, Statuscodes sowie die
Orchestrierung einzelner Use Cases. Datenbankzugriffe sind in
dedizierten Repositories gekapselt. Dadurch bleibt die API-Schicht frei
von SQL-Details, während Persistenzlogik zentral gebündelt und
wiederverwendbar gehalten wird.

Für LLM-nahe Abläufe wird ergänzend eine Service-Schicht genutzt
(insbesondere \passthrough{\lstinline!ticker\_service!} und
\passthrough{\lstinline!llm\_service!}). Die Service-Schicht wird dabei
gezielt, aber nicht flächendeckend eingesetzt: Viele Standard-CRUD- und
Listenoperationen folgen dem Muster Router → Repository, während
komplexere Generierungs- und Orchestrierungspfade über Services laufen.

Für die Datenrepräsentation gilt eine bewusste Trennung: ORM-Modelle
definieren die Persistenzstruktur, Pydantic-Schemas die API-Verträge.
Das stabilisiert die Schnittstellen gegenüber internen Refactorings und
verbessert Wartbarkeit und Testbarkeit.

\begin{center}\rule{0.5\linewidth}{0.5pt}\end{center}

\subsubsection{4.2.3 API-Design und
Endpunktstrategie}\label{api-design-und-endpunktstrategie}

Die Backend-Schnittstelle ist als versionierte REST-API unter
\passthrough{\lstinline!/api/v1!} ausgelegt und orientiert sich am
Architekturstil \textbf{Representational State Transfer} (Fielding 2000,
Kap. 5). Endpunkte sind ressourcenorientiert strukturiert (z. B.
\passthrough{\lstinline!teams!}, \passthrough{\lstinline!matches!},
\passthrough{\lstinline!ticker!}, \passthrough{\lstinline!media!},
\passthrough{\lstinline!clips!}) und folgen konsistenten Benennungs- und
Methodenregeln. \passthrough{\lstinline!GET!} dient lesenden Abfragen,
\passthrough{\lstinline!POST!} dem Erzeugen bzw. Triggern,
\passthrough{\lstinline!PATCH!} partiellen Zustandsänderungen und
\passthrough{\lstinline!DELETE!} kontrollierten Löschoperationen.
Insbesondere im Ticker-Kontext bildet \passthrough{\lstinline!PATCH!}
den redaktionellen Lebenszyklus (\passthrough{\lstinline!draft!} →
\passthrough{\lstinline!published!} bzw.
\passthrough{\lstinline!rejected!}) explizit ab.

Konzeptionell wird zwischen fachlichen Datenendpunkten und prozessualen
Triggerendpunkten getrennt. Dadurch bleiben Leseoperationen
seiteneffektfrei, während Import- und Generierungsprozesse gezielt
ausgelöst werden können. Fehlerfälle werden konsistent über
HTTP-Statuscodes und strukturierte Fehlermeldungen kommuniziert, sodass
Frontend und n8n-Workflows deterministisch reagieren können.

\begin{center}\rule{0.5\linewidth}{0.5pt}\end{center}

\subsubsection{4.2.4 Asynchronität, Nebenläufigkeit und
Performance}\label{asynchronituxe4t-nebenluxe4ufigkeit-und-performance}

Die Laufzeitkonzeption setzt Asynchronität gezielt auf I/O-intensiven
Pfaden ein, insbesondere bei LLM-bezogenen Generierungsrouten,
Media-Verarbeitung und WebSocket-Kommunikation. Das erhöht die
Reaktionsfähigkeit unter paralleler Last, da Wartezeiten auf externe
Dienste keine Worker dauerhaft blockieren.

Die Anzahl gleichzeitiger LLM-Aufrufe wird über einen konfigurierbaren
\textbf{Semaphore-Mechanismus} begrenzt (Standard: 8 parallele
Anfragen). Diese Begrenzung reduziert Lastspitzen und verringert die
Wahrscheinlichkeit von Rate-Limit-Problemen seitens der LLM-Provider.
Der Semaphore wirkt dabei pro Prozessinstanz --- bei horizontaler
Skalierung steigt die effektive Gesamtkonkurrenz entsprechend der Anzahl
laufender Instanzen. Ergänzend greift ein \textbf{Retry-Mechanismus mit
Backoff}: Bei Rate-Limit-Fehlern erfolgen automatische
Wiederholungsversuche mit linearem Backoff (3 Versuche, 30 s / 60 s / 90
s); bei sonstigen transienten Fehlern exponentielles Backoff (1 s / 2 s
/ 4 s).

Die Implementierung ist bewusst hybrid aus synchronen und asynchronen
Pfaden aufgebaut. Nicht alle Endpunkte sind vollständig asynchronisiert;
stattdessen wurde ein pragmatischer Kompromiss zwischen Performance,
Komplexität und Wartbarkeit gewählt. Langlaufende Integrationslogik ist
in n8n ausgelagert, wodurch der API-Server als transaktionaler Kern
entlastet wird.

\begin{center}\rule{0.5\linewidth}{0.5pt}\end{center}

\subsubsection{4.2.5 Fehlerbehandlung, Robustheit und
Betriebsaspekte}\label{fehlerbehandlung-robustheit-und-betriebsaspekte}

Für den Live-Betrieb wurde das Backend auf robuste Fehlerbehandlung und
kontrollierte Degradation ausgelegt. Fehler werden über konsistente
HTTP-Antworten und strukturierte Meldungen zurückgegeben, sodass
Frontend und n8n-Workflows differenziert reagieren können.

Die Datenpersistenz ist auf \textbf{idempotente Verarbeitung} ausgelegt
(vgl. Kap. 4.3.4 für Details zu Upsert-Strategien und Schlüsseldesign).
Bei Teilausfällen externer Dienste bleibt das System eingeschränkt
funktionsfähig --- bereits persistierte Daten bleiben verfügbar,
redaktionelle Kernprozesse (z. B. manuelle Tickererstellung) können
fortgeführt werden.

\begin{center}\rule{0.5\linewidth}{0.5pt}\end{center}

\subsubsection{4.2.6 Sicherheitskonzept}\label{sicherheitskonzept}

Das System adressiert Sicherheit auf drei Ebenen, wobei der
Projektrahmen einer Bachelorarbeit eine bewusste Priorisierung
erfordert:

\textbf{Transportebene:} Die gesamte Kommunikation zwischen Frontend,
Backend und externen Diensten erfolgt über HTTPS. Die Render-Plattform
übernimmt die TLS-Terminierung und Zertifikatsverwaltung automatisiert.

\textbf{API-Absicherung:} Der Backend-Server konfiguriert eine
\textbf{CORS-Whitelist} (Cross-Origin Resource Sharing), die Anfragen
ausschließlich von autorisierten Frontend-Ursprüngen akzeptiert. Die
API-Schlüssel externer Dienste (LLM-Provider, API-Football, ScorePlay)
werden über Umgebungsvariablen injiziert und sind weder im Quellcode
noch im Frontend exponiert. Die Pydantic-basierte Eingabevalidierung
schützt zusätzlich vor Injection-Angriffen, da sämtliche
Request-Payloads gegen typisierte Schemas geprüft werden, bevor sie die
Geschäftslogik erreichen.

\textbf{Authentifizierung und Autorisierung:} Eine feingranulare
Benutzerauthentifizierung mit Rollenkonzept (z. B. Redakteur,
Administrator) ist im aktuellen Stand \textbf{nicht implementiert}. Das
System ist als internes Redaktionswerkzeug konzipiert und setzt auf
Netzwerksegmentierung statt individueller Zugangskontrolle. Für einen
produktiven Mehrbenutzerbetrieb wäre die Integration eines
Authentifizierungsframeworks (z. B. OAuth 2.0 oder JWT-basierte
Token-Authentifizierung) erforderlich --- dies wird in Kapitel 7 als
Erweiterungsperspektive eingeordnet.

\begin{center}\rule{0.5\linewidth}{0.5pt}\end{center}

\subsection{4.3 Datenbankkonzeption}\label{datenbankkonzeption}

Die Persistenzschicht basiert auf PostgreSQL und umfasst in der
aktuellen Fassung \textbf{18 Tabellen} (17 ORM-Modelle sowie eine
schlüsselwertbasierte \passthrough{\lstinline!settings!}-Tabelle, die
ausschließlich über eine Alembic-Migration verwaltet wird). Die Wahl
eines relationalen Systems gegenüber dokumentenorientierten Ansätzen (z.
B. MongoDB) begründet sich durch die stark strukturierten Domänendaten:
FK-Beziehungen zwischen Teams, Wettbewerben, Spielen, Events und
Ticker-Einträgen sind klar definiert und referenziell integer zu halten.
JSONB-Felder werden selektiv für semistrukturierte Daten (z. B.
Statistik-Rohdaten) eingesetzt, sodass die ACID-Garantien relationaler
Datenbanken erhalten bleiben. Das Schema ist auf einen stabilen
Live-Betrieb mit wiederholbaren Importen, klaren Zustandsübergängen und
nachvollziehbaren Redaktionsentscheidungen ausgelegt.

\subsubsection{4.3.1
Schemadesign-Prinzipien}\label{schemadesign-prinzipien}

Die folgende Abbildung zeigt die zentralen Entitäten und ihre
Beziehungen:

\begin{lstlisting}
erDiagram
    countries ||--o{ teams : 'hat Teams'
    teams ||--o{ competition_teams : 'spielt in'
    competitions ||--o{ competition_teams : 'hat Teams'
    competitions ||--o{ matches : 'umfasst'
    seasons ||--o{ matches : 'enthält'
    teams ||--o{ matches : 'home_team'
    teams ||--o{ matches : 'away_team'

    matches ||--o{ events : 'hat Events'
    matches ||--o{ synthetic_events : 'hat synth. Events'
    matches ||--o{ ticker_entries : 'hat Ticker'
    matches ||--o{ lineups : 'hat Lineup'
    matches ||--o{ match_statistics : 'hat Statistiken'

    events ||--o{ ticker_entries : 'event_id (nullable)'
    synthetic_events |o--o{ ticker_entries : 'synthetic_event_id (nullable)'

    players ||--o{ lineups : 'spielt'
    players ||--o{ player_statistics : 'hat Statistiken'

    seasons ||--o{ standings : 'hat Standings'
    competitions ||--o{ standings : 'hat Standings'
    teams ||--o{ standings : 'Rang'
\end{lstlisting}

\begin{quote}
\textbf{Hinweis:} Die Tabellen \passthrough{\lstinline!media\_queue!},
\passthrough{\lstinline!media\_clips!},
\passthrough{\lstinline!style\_references!} und
\passthrough{\lstinline!settings!} sind als eigenständige Entitäten ohne
Fremdschlüsselbeziehungen modelliert und daher im ER-Diagramm nicht
abgebildet. \passthrough{\lstinline!style\_references!} dient als
Few-Shot-Datenquelle für die LLM-Promptgenerierung;
\passthrough{\lstinline!media\_queue!} wird applikationsseitig über die
Media-Endpunkte verwaltet.
\end{quote}

Das Schema folgt einem hybriden Entwurfsansatz aus strukturierter
Normalisierung und gezielter Schemaflexibilität:

\begin{enumerate}
\def\labelenumi{\arabic{enumi}.}

\item
  \textbf{Normalisierte Kernentitäten} --- Zentrale Domänenobjekte wie
  \passthrough{\lstinline!teams!}, \passthrough{\lstinline!matches!},
  \passthrough{\lstinline!events!},
  \passthrough{\lstinline!ticker\_entries!},
  \passthrough{\lstinline!players!},
  \passthrough{\lstinline!competitions!} und Zuordnungstabellen sind in
  dritter Normalform modelliert (Codd 1970). Dadurch bleiben
  Beziehungen, Integrität und Auswertbarkeit im Live-Betrieb stabil.
\item
  \textbf{Gezielter JSONB-Einsatz} --- JSONB wird bewusst nur dort
  eingesetzt, wo sich Strukturen dynamisch ändern oder kontextabhängig
  sind (z. B. \passthrough{\lstinline!synthetic\_events.data!}, sowie
  mehrere Kontextfelder in \passthrough{\lstinline!matches!}). Die
  Statistiktabellen (\passthrough{\lstinline!match\_statistics!},
  \passthrough{\lstinline!player\_statistics!}) sind dagegen bewusst
  stark typisiert über explizite Spalten modelliert, um konsistente
  Abfragen und robuste Aggregationen zu ermöglichen.
\item
  \textbf{Schlüsselstrategie: Integer-PK + UUID-Fachschlüssel} --- Das
  Schema verwendet numerische Primärschlüssel
  (\passthrough{\lstinline!id!}) für performante Joins und
  Fremdschlüsselbeziehungen. Zusätzlich existiert in mehreren
  Kernentitäten ein eindeutiges \passthrough{\lstinline!uid!}-Feld
  (UUID) als stabiler externer Referenzschlüssel. Es handelt sich damit
  nicht um ein reines UUID-Primary-Key-Schema, sondern um eine
  kombinierte Strategie.
\end{enumerate}

\begin{center}\rule{0.5\linewidth}{0.5pt}\end{center}

\subsubsection{4.3.2 Status-Lifecycle der
Ticker-Einträge}\label{status-lifecycle-der-ticker-eintruxe4ge}

Jeder Ticker-Eintrag in \passthrough{\lstinline!ticker\_entries!} folgt
einem klaren Statusmodell mit drei Zuständen:

\begin{itemize}

\item
  \textbf{\passthrough{\lstinline!draft!}}: Entwurf, noch nicht
  freigegeben
\item
  \textbf{\passthrough{\lstinline!published!}}: redaktionell oder
  automatisch veröffentlicht
\item
  \textbf{\passthrough{\lstinline!rejected!}}: verworfen, bleibt zur
  Nachvollziehbarkeit erhalten
\end{itemize}

\begin{lstlisting}
stateDiagram-v2
    [*] --> draft : KI generiert (coop)\nauto_publish=false
    [*] --> published : KI generiert (auto)\nauto_publish=true
    [*] --> published : Manuell\nPOST /ticker/manual

    draft --> published : Redakteur bestätigt\nPATCH status=published
    draft --> rejected  : Redakteur verwirft\nPATCH status=rejected
    published --> draft : PATCH status=draft

    note right of rejected : Bleibt erhalten\n(Auswertbarkeit)
\end{lstlisting}

Eine Rückstufung von \passthrough{\lstinline!published!} auf
\passthrough{\lstinline!draft!} ist über
\passthrough{\lstinline!PATCH status=draft!} möglich und wird
insbesondere bei Re-Triggern von Matchphasen-Workflows genutzt (vgl.
Kap. 5.5.3). \passthrough{\lstinline!rejected!} ist ein terminaler
Zustand --- ein verworfener Eintrag kann nicht reaktiviert werden.

Ergänzend markiert das Feld \passthrough{\lstinline!source!} die
Herkunft (\passthrough{\lstinline!ai!}
vs.~\passthrough{\lstinline!manual!}) und erlaubt eine saubere Trennung
für Evaluationen (z. B. ausschließlich KI-generierte Einträge).

\begin{quote}
\textbf{Hinweis für die Methodik}: Im aktuellen Schema ist kein
separates \passthrough{\lstinline!published\_at!}-Feld vorhanden.
Persistiert wird \passthrough{\lstinline!created\_at!} des jeweiligen
Ticker-Eintrags. Für sekundengenaue Time-to-Publish-Analysen sind daher
zusätzliche Messzeitpunkte im Evaluationsdatensatz bzw. in der Messlogik
erforderlich, statt sich allein auf den Tabellenzustand zu stützen.
\end{quote}

\begin{center}\rule{0.5\linewidth}{0.5pt}\end{center}

\subsubsection{4.3.3 Die drei Betriebsmodi}\label{die-drei-betriebsmodi}

Das Feld \passthrough{\lstinline!ticker\_mode!} in
\passthrough{\lstinline!matches!} steuert das Verhalten der
Generierungspipeline pro Spiel und ist zur Laufzeit umschaltbar über
\passthrough{\lstinline!PATCH /api/v1/matches/\{id\}/ticker-mode!}. Es
sind drei Modi implementiert:

\begin{enumerate}
\def\labelenumi{\arabic{enumi}.}

\item
  \textbf{\passthrough{\lstinline!auto!}} (vollautomatisch) ---
  Triggerkette läuft ohne redaktionellen Eingriff; KI-Ergebnisse werden
  direkt veröffentlicht.
\item
  \textbf{\passthrough{\lstinline!coop!}} (kooperativ /
  Human-in-the-Loop) --- KI erzeugt Entwürfe
  (\passthrough{\lstinline!draft!}), die Redaktion entscheidet über
  Veröffentlichung oder Ablehnung. Dieser Modus bildet das primäre
  Zielbild des Systems.
\item
  \textbf{\passthrough{\lstinline!manual!}} (rein redaktionell) ---
  Einträge werden manuell erstellt; KI-Generierung ist nicht leitend für
  den Veröffentlichungsprozess. Im
  \passthrough{\lstinline!manual!}-Modus werden keine
  KI-Generierungsendpunkte getriggert; eingehende Events lösen
  ausschließlich die Datenpersistenz aus, nicht die LLM-Pipeline. Dieser
  Modus dient als Referenz für Vergleiche in der Evaluation.
\end{enumerate}

Der Moduswechsel ist als Laufzeitparameter auf Datenbankebene
abgebildet; der Statusfluss ist im State-Diagramm in Abschnitt 4.3.2
dargestellt.

\begin{center}\rule{0.5\linewidth}{0.5pt}\end{center}

\subsubsection{4.3.4 Integrität, Idempotenz und
Auswertbarkeit}\label{integrituxe4t-idempotenz-und-auswertbarkeit}

Für den Zusammenschluss aus n8n-Workflows, externen APIs und
Backend-Routen ist die Datenbank auf idempotente Verarbeitung ausgelegt:

\begin{enumerate}
\def\labelenumi{\arabic{enumi}.}

\item
  \textbf{Konfliktarme Wiederholbarkeit} --- Wiederholte Imports führen
  durch Unique-Constraints und Upsert-Strategien nicht zu
  unkontrollierten Duplikaten.
\item
  \textbf{Eindeutige Zuordnung in zentralen Tabellen} --- Eindeutige
  Schlüssel (z. B. \passthrough{\lstinline!source\_id!} pro
  Ereignisquelle) sichern die Nachvollziehbarkeit importierter
  Ereignisse.
\item
  \textbf{Referenzielle Integrität über FK-Regeln} --- Lösch- und
  Nullsetz-Strategien (\passthrough{\lstinline!CASCADE!},
  \passthrough{\lstinline!SET NULL!}) sorgen dafür, dass abhängige Daten
  konsistent bleiben und Historie dort erhalten wird, wo sie fachlich
  relevant ist.
\item
  \textbf{Auswertungskompatibilität} --- Die Kombination aus
  \passthrough{\lstinline!ticker\_mode!},
  \passthrough{\lstinline!status!}, \passthrough{\lstinline!source!},
  Spielbezug und Zeitfeldern schafft die Grundlage für reproduzierbare
  Vergleichsauswertungen zwischen manuellen, kooperativen und
  automatisierten Abläufen.
\end{enumerate}

\begin{center}\rule{0.5\linewidth}{0.5pt}\end{center}

\subsection{4.4 Workflow-Konzeption
(n8n)}\label{workflow-konzeption-n8n}

Die Workflow-Schicht ist als entkoppelte Orchestrierungsebene zwischen
externen Datenquellen und Backend ausgelegt. n8n wurde gegenüber
Alternativen wie einem Custom-Scheduler (z. B. Celery, APScheduler) oder
cron-basierten Skripten gewählt, weil es Workflow-Logik ohne
Codeänderungen am Backend anpassbar macht, eine visuelle
Debugging-Oberfläche bietet und HTTP-Trigger sowie Webhook-Empfang ohne
zusätzliche Infrastruktur unterstützt. n8n übernimmt dabei API-Aufrufe,
Transformation, Persistenzvorbereitung und Triggersteuerung, während das
FastAPI-Backend als transaktionaler Kern und Integrationspunkt für
Frontend und KI-Generierung fungiert.

\subsubsection{4.4.1 Workflow-Klassen und
Verantwortlichkeiten}\label{workflow-klassen-und-verantwortlichkeiten}

Alle Generierungs-Workflows folgen einem konsistenten
Vier-Phasen-Muster: \textbf{(1) Webhook-Trigger} --- ein HTTP-Aufruf
löst den Workflow aus; \textbf{(2) SQL-Operation} --- ein
Upsert-Statement persistiert die Daten mit Konfliktbehandlung
(\passthrough{\lstinline!ON CONFLICT!}); \textbf{(3) Filter-Knoten} ---
nur neue, nicht bereits vorhandene Zeilen (erkennbar an der
zurückgegebenen Datenbank-\passthrough{\lstinline!id!}) passieren
weiter; \textbf{(4) Backend-Call} --- der entsprechende FastAPI-Endpunkt
wird pro neuer Zeile aufgerufen, bei Generierungsworkflows zum
LLM-Trigger. Dieses Muster garantiert Idempotenz: Mehrfachaufrufe
desselben Events erzeugen weder Datenduplikate noch redundante
LLM-Aufrufe.

Die n8n-Landschaft gliedert sich in vier funktionale Klassen. Die
folgende Übersicht zeigt alle 15 Workflows und ihren Auslöser:

{\def\LTcaptype{none} % do not increment counter
\begin{longtable}[]{@{}
  >{\raggedright\arraybackslash}p{(\linewidth - 4\tabcolsep) * \real{0.2727}}
  >{\raggedright\arraybackslash}p{(\linewidth - 4\tabcolsep) * \real{0.4091}}
  >{\raggedright\arraybackslash}p{(\linewidth - 4\tabcolsep) * \real{0.3182}}@{}}
\toprule\noalign{}
\begin{minipage}[b]{\linewidth}\raggedright
Gruppe
\end{minipage} & \begin{minipage}[b]{\linewidth}\raggedright
Workflows
\end{minipage} & \begin{minipage}[b]{\linewidth}\raggedright
Trigger
\end{minipage} \\
\midrule\noalign{}
\endhead
\bottomrule\noalign{}
\endlastfoot
\textbf{A) Stammdaten} & 01--04: Countries, Teams, Competitions, Matches
& Bei leeren Daten \\
\textbf{B) Matchdaten} & 05--07: Lineups, Match-/Player-Statistics,
Prematch & Nach Matchauswahl \\
\textbf{C) KI-Generierung} & 09 Events-LLM, 13 Halftime/Aftertime, 14
Anpfiff/Abpfiff & Modus- und Event-gesteuert \\
\textbf{D) Medien \& Social} & 08 ScorePlay, 10 Twitter/X, 11 YouTube,
12 Instagram & Medien-Suche \\
\end{longtable}
}

Diese Trennung reduziert Kopplung, erleichtert Fehlersuche und erlaubt
es, einzelne Teilprozesse unabhängig anzupassen. Die konkreten
Workflow-Dateinamen und Implementierungsdetails sind in Kapitel 5.5
dokumentiert.

\begin{center}\rule{0.5\linewidth}{0.5pt}\end{center}

\subsubsection{4.4.2 Workflow-Grenzen und projektspezifische
Besonderheiten}\label{workflow-grenzen-und-projektspezifische-besonderheiten}

Die Workflow-Landschaft ist funktional umfassend, enthält aber bewusst
pragmatische Projektentscheidungen:

\begin{enumerate}
\def\labelenumi{\arabic{enumi}.}

\item
  \textbf{Teilweise saisonspezifische Fixierung} --- In einzelnen Flows
  ist die Saisonlogik auf 2025 begrenzt.
\item
  \textbf{Hybride Aktualisierungslogik} --- Nicht alle Datenpfade sind
  eventgetrieben; ein Teil wird durch Polling und bedarfsgesteuerte
  Trigger ergänzt.
\item
  \textbf{Instanzspezifische Spezialpfade} --- Einige
  LLM-/Medienprozesse enthalten EF-spezifische Logik, was für eine
  vollständig generalisierte White-Label-Fähigkeit parametriert werden
  müsste.
\end{enumerate}

In Summe bildet n8n eine tragfähige Orchestrierungsschicht, die externe
Datenquellen, interne Persistenz und KI-generierte Inhalte in einem
modularen, nachvollziehbaren und erweiterbaren Ablauf zusammenführt.

\begin{center}\rule{0.5\linewidth}{0.5pt}\end{center}

\subsection{4.5 KI-Komponente}\label{ki-komponente}

\subsubsection{4.5.1
Multi-Provider-Architektur}\label{multi-provider-architektur}

Die KI-Komponente ist als providerunabhängige Abstraktionsschicht
implementiert. Der LLM-Dienst kapselt mehrere Anbieter hinter einer
einheitlichen Schnittstelle und unterstützt aktuell OpenAI, Anthropic,
Google Gemini, OpenRouter sowie einen Mock-Modus für Entwicklungs- und
Fallbackszenarien ohne API-Key.

\textbf{Primärzugang: OpenRouter als API-Gateway}

Das vorliegende System setzt primär \textbf{OpenRouter} als LLM-Gateway
ein. OpenRouter fungiert als einheitliche API-Schnittstelle, über die
mit einem einzelnen API-Schlüssel verschiedene Sprachmodelle
unterschiedlicher Anbieter angesprochen werden können. Diese
Architekturentscheidung bietet drei Vorteile:

\begin{enumerate}
\def\labelenumi{\arabic{enumi}.}

\item
  \textbf{Kein Vendor Lock-in}: Der Wechsel zwischen Modellen erfordert
  lediglich eine Änderung der Umgebungsvariable
  \passthrough{\lstinline!OPENROUTER\_MODEL!}, keine Codeanpassung.
\item
  \textbf{Vergleichbarkeit}: Verschiedene Modelle können auf demselben
  Prompt evaluiert werden (vgl. Abschnitt 6.3.3).
\item
  \textbf{Kostenoptimierung}: Das jeweils kostengünstigste Modell für
  die Aufgabe kann gewählt werden.
\end{enumerate}

Im produktiven Einsatz wird aktuell \textbf{Gemini 2.0 Flash Lite}
(Modell-ID \passthrough{\lstinline!google/gemini-2.0-flash-lite-001!})
über OpenRouter genutzt. Die Wahl dieses kompakten Modells begründet
sich durch die Aufgabencharakteristik (vgl. Kap. 3.1): Liveticker-Texte
sind kurz und faktenbasiert --- eine Aufgabe, für die kompakte Modelle
eine vergleichbare Qualität bei deutlich niedrigerer Latenz erzielen.
Für Vereinsredaktionen mit begrenztem Budget ist dieser Kostenvorteil
ein entscheidender Faktor für die Praxistauglichkeit des Systems.

\textbf{Fallback-Kette und Mock-Provider}

Ergänzend implementiert das Backend eine \textbf{Fallback-Kette} für den
Fall, dass OpenRouter nicht konfiguriert ist. Das System prüft zur
Startzeit in der festen Prioritätsreihenfolge, ob ein API-Schlüssel
hinterlegt ist, und aktiviert den ersten verfügbaren Provider als
Singleton:

\begin{quote}
\passthrough{\lstinline!openrouter!} → \passthrough{\lstinline!gemini!}
→ \passthrough{\lstinline!openai!} → \passthrough{\lstinline!anthropic!}
→ \passthrough{\lstinline!mock!}
\end{quote}

Ist kein Schlüssel konfiguriert, wird auf einen regelbasierten
\textbf{Mock-Provider} zurückgefallen, der Template-basierte Texte ohne
LLM-Aufruf erzeugt. Diese Architektur stellt sicher, dass das System
auch ohne LLM-Anbindung lauffähig bleibt --- etwa für Entwicklung und
Tests.

Standardmodelle der Provider sind:

{\def\LTcaptype{none} % do not increment counter
\begin{longtable}[]{@{}ll@{}}
\toprule\noalign{}
Provider & Standardmodell \\
\midrule\noalign{}
\endhead
\bottomrule\noalign{}
\endlastfoot
OpenRouter &
\passthrough{\lstinline!google/gemini-2.0-flash-lite-001!} \\
Gemini & \passthrough{\lstinline!gemini-2.0-flash-lite-001!} \\
OpenAI & \passthrough{\lstinline!gpt-4o-mini!} \\
Anthropic & \passthrough{\lstinline!claude-haiku-4-5-20251001!} \\
\end{longtable}
}

Zusätzlich kann die Auswahl in Generierungsrouten pro Request über
\passthrough{\lstinline!provider!} und \passthrough{\lstinline!model!}
überschrieben werden, was insbesondere für Evaluationsvergleiche
zwischen Modellen genutzt wird (vgl. Abschnitt 6.3.3).

\begin{center}\rule{0.5\linewidth}{0.5pt}\end{center}

\subsubsection{4.5.2
Prompt-Engineering-Strategie}\label{prompt-engineering-strategie}

Die Generierung verwendet ein template-basiertes Prompting mit modularen
Bausteinen, die dynamisch zusammengesetzt werden:

\begin{enumerate}
\def\labelenumi{\arabic{enumi}.}

\item
  Rolleninstruktion (Persona, Aufgabe, Gattung)
\item
  Stilinstruktion (neutral / euphorisch / kritisch)
\item
  Faktenblock zum Ereignis (Typ, Detail, Minute, beteiligte
  Spieler/Team)
\item
  Dynamischer Kontextblock mit relevanten Matchinformationen
\item
  Optionaler Few-Shot-Block mit Stilbeispielen
\item
  Regelblock mit formativen und inhaltlichen Einschränkungen
\end{enumerate}

\begin{lstlisting}
flowchart TD
    A['Ereignis-Typ\n(z.B. goal)'] --> P
    B['Fakten\n(Spieler, Minute, Team)'] --> P
    C['Match-Kontext\n(Spielstand, Teams, Liga)'] --> P
    D['Few-Shot-Beispiele\naus style_references\n(0–3 Texte)'] --> P
    E['Regelblock\n(Länge, Sprache, kein Markdown)'] --> P
    F['Stilinstruktion\nneutral / euphorisch / kritisch'] --> P

    P['Vollständiger\nSystem-Prompt'] --> LLM['LLM\n(temp=0.3)']
    LLM --> T['Generierter\nTicker-Text']
\end{lstlisting}

Der Few-Shot-Block wird aus der Tabelle
\passthrough{\lstinline!style\_references!} gespeist, die manuell
kuratierte Original-Tickertexte von Eintracht Frankfurt enthält. Pro
LLM-Aufruf werden bis zu drei zufällige Referenzen selektiert, gefiltert
nach \passthrough{\lstinline!event\_type!},
\passthrough{\lstinline!instance!} und optional
\passthrough{\lstinline!league!}. Durch die Randomisierung wird eine
monotone Reproduktion vermieden, während der stilistische Korridor
gewahrt bleibt. Für Pre-Match-Typen enthält der Prompt zusätzliche harte
Restriktionen, um Live-Szenen-Halluzinationen zu vermeiden und die
Ausgabe auf Vorschau- und Analyseinhalte zu begrenzen.

Der Inferenzparameter \textbf{Temperature} wird für die Textgenerierung
auf \passthrough{\lstinline!0.3!} festgelegt, um die faktische
Korrektheit gegenüber kreativer Varianz zu priorisieren. Für
Übersetzungsaufgaben wird ein noch niedrigerer Wert von
\passthrough{\lstinline!0.1!} verwendet, um semantische Abweichungen vom
Originaltext zu minimieren.

\begin{center}\rule{0.5\linewidth}{0.5pt}\end{center}

\subsubsection{4.5.3 Stilprofile und
Instanzspezifik}\label{stilprofile-und-instanzspezifik}

Das System unterstützt drei zentrale Stilprofile:

\begin{itemize}

\item
  \textbf{\passthrough{\lstinline!neutral!}}: sachlich, ausgewogen, ohne
  Vereinspräferenz
\item
  \textbf{\passthrough{\lstinline!euphorisch!}}: emotional und fan-nah
\item
  \textbf{\passthrough{\lstinline!kritisch!}}: analytisch und bewertend
\end{itemize}

Ergänzend steuert die Instanzkonfiguration
(\passthrough{\lstinline!generic!}
vs.~\passthrough{\lstinline!ef\_whitelabel!}) die Tonalität. Für
\passthrough{\lstinline!ef\_whitelabel!} werden vereinsspezifische
Stilbeispiele aus \passthrough{\lstinline!style\_references!} als
Few-Shot-Kontext eingebunden; die generische Instanz kann ohne
vereinsgebundene Stilprägung betrieben werden.

\begin{center}\rule{0.5\linewidth}{0.5pt}\end{center}

\subsubsection{4.5.4 Mehrsprachigkeit}\label{mehrsprachigkeit}

Die Mehrsprachigkeit ist als First-Class-Parameter der LLM-Pipeline
implementiert: Die Zielsprache (\passthrough{\lstinline!language!},
Standard \passthrough{\lstinline!de!}) wird direkt an den Prompt
übergeben, sodass der Text in der gewünschten Sprache generiert wird ---
nicht nachträglich übersetzt. Dieses Vorgehen nutzt die Fähigkeit
moderner großer Sprachmodelle, idiomatische Texte ohne den
Zwischenschritt einer maschinellen Übersetzung zu produzieren (vgl. Kap.
3.1). Die Generierungsendpunkte akzeptieren jeden ISO-639-1-Sprachcode;
getestet sind Deutsch (\passthrough{\lstinline!de!}), Englisch
(\passthrough{\lstinline!en!}) und Japanisch
(\passthrough{\lstinline!ja!}) --- die drei in Kapitel 2.2 als relevant
identifizierten Sprachen.

Für bestehende Einträge steht eine Batch-Übersetzung über
\passthrough{\lstinline!POST /api/v1/ticker/translate-batch/\{match\_id\}!}
zur Verfügung. Sie übersetzt alle veröffentlichten KI-Einträge eines
Spiels mit reduzierter Temperatur (\passthrough{\lstinline!0.1!}), um
möglichst deterministische Ergebnisse zu erzielen. Der aktuelle
Implementierungsstand schließt dabei auch manuell erstellte Einträge ein
--- eine Einschränkung, die in Kapitel 6.7 dokumentiert ist.

Eine verbleibende Limitation betrifft die Few-Shot-Referenzen: Die
\passthrough{\lstinline!style\_references!}-Tabelle enthält
ausschließlich deutschsprachige Beispieltexte. Bei nicht-deutschen
Generierungen fehlt damit die stilprägende Few-Shot-Komponente, was die
Qualität der Tonalität gegenüber dem deutschen Modus leicht abschwächen
kann. Eine Erweiterung der Stildatenbank um englisch- und
japanischsprachige Referenzen wäre der naheliegende Folgeschritt (vgl.
Kap. 8.3.2).

\begin{center}\rule{0.5\linewidth}{0.5pt}\end{center}

\subsection{4.6 Frontend-Konzeption}\label{frontend-konzeption}

\subsubsection{4.6.1 Architekturprinzipien}\label{architekturprinzipien}

Das Frontend ist als React-Client mit TypeScript und klarer Trennung
zwischen UI-Komponenten, datenbezogenen Hooks und API-Zugriffsschicht
umgesetzt. Der zentrale Anwendungszustand wird überwiegend lokal bzw.
feature-nah gehalten; gemeinsamer Zustand für tief verschachtelte
Komponenten wird über drei dedizierte Contexts bereitgestellt
(\passthrough{\lstinline!TickerModeContext!},
\passthrough{\lstinline!TickerDataContext!},
\passthrough{\lstinline!TickerActionsContext!}), um Prop-Drilling zu
vermeiden und Re-Renders besser zu kontrollieren.

Die White-Label-Fähigkeit ist über eine Konfigurationsschicht
(\passthrough{\lstinline!config/whitelabel.ts!}) realisiert (u. a.
Farben, Team-Keyword, API-Base-URLs). Die Oberfläche unterstützt alle
drei Betriebsmodi (\passthrough{\lstinline!auto!},
\passthrough{\lstinline!coop!}, \passthrough{\lstinline!manual!})
inklusive Laufzeitumschaltung über den
\passthrough{\lstinline!ModeSelector!} und Synchronisierung mit dem
Backend.

\begin{center}\rule{0.5\linewidth}{0.5pt}\end{center}

\subsubsection{4.6.2 Hook- und
State-Architektur}\label{hook--und-state-architektur}

Die Zustands- und Interaktionslogik ist in wiederverwendbare Custom
Hooks aufgeteilt --- ein Entwurfsmuster, das React seit Version 16.8 als
Alternative zu Klassenkomponenten eingeführt hat. Das Frontend umfasst
25 spezialisierte Hooks, die sich in fünf funktionale Kategorien
gliedern lassen:

\begin{enumerate}
\def\labelenumi{\arabic{enumi}.}

\item
  \textbf{Navigation und Datenauswahl} --- Steuerung der hierarchischen
  Auswahlkette (Land → Team → Wettbewerb → Spieltag → Spiel) mit
  bedarfsgesteuerter Auslösung von n8n-Importtriggern bei leerem
  Datenbestand.
\item
  \textbf{Datenabfrage und Polling} --- Periodisches Laden von Match-,
  Event- und Tickerdaten über drei spezialisierte Hooks, die jeweils
  einen eigenen Abfragezyklus betreiben und über einen aggregierenden
  Hook koordiniert werden.
\item
  \textbf{Echtzeitkommunikation} --- WebSocket-Verbindung für den
  Medien-Feed mit automatischer Reconnect-Strategie.
\item
  \textbf{Redaktionsinteraktion} --- Modusverwaltung,
  Draft-Bestätigung/-Ablehnung, Tastatursteuerung sowie automatische
  Publikation im Vollautomatik-Modus.
\item
  \textbf{UI-Steuerung} --- Panel-Größenanpassung,
  Spielminuten-Berechnung, Spielernamen-Autovervollständigung und
  Systemstatus-Überwachung.
\end{enumerate}

\begin{center}\rule{0.5\linewidth}{0.5pt}\end{center}

\subsubsection{4.6.3 Dreispalten-Layout und
Redaktionsfluss}\label{dreispalten-layout-und-redaktionsfluss}

Die Hauptansicht folgt einem responsiven Dreispalten-Ansatz mit klarer
Aufgabenverteilung:

\begin{enumerate}
\def\labelenumi{\arabic{enumi}.}

\item
  \textbf{Linkes Panel} --- Veröffentlichungsperspektive: listet
  publizierte Ticker-Einträge als chronologischen Feed, dedupliziert
  nach \passthrough{\lstinline!event\_id!} und sortiert nach
  Spielminute.
\item
  \textbf{Mittleres Panel} --- Redaktionelle Kernarbeit:
  Event-Verarbeitung in Abhängigkeit vom Modus (auto, coop, manual),
  Draft-Prüfung und Freigabe/Ablehnung, manuelle Eingabe mit
  Command-Parser (z. B. \passthrough{\lstinline!/g!},
  \passthrough{\lstinline!/gelb!}, \passthrough{\lstinline!/rot!},
  \passthrough{\lstinline!/s!}), sowie am unteren Rand integrierte
  Media- und Social-Panels (ScorePlay-Bilder, YouTube, Twitter/X,
  Instagram). Die Social-Media-Panels werden ausschließlich bei
  erkannten EF-Spielen (\passthrough{\lstinline!isOurTeam!})
  eingeblendet.
\item
  \textbf{Rechtes Panel} --- Kontext- und Analyseinformationen:
  Matchstatistiken, Aufstellungen inkl. Wechsel-/Karten-/Injury-Kontext
  sowie spielerbezogene Leistungsdaten.
\end{enumerate}

\begin{center}\rule{0.5\linewidth}{0.5pt}\end{center}

\subsubsection{4.6.4 Moduslogik und
Interaktionsdesign}\label{moduslogik-und-interaktionsdesign}

Die Modusumschaltung (vgl. Kap. 4.3.3 für Modi-Definition) ist zentraler
Bestandteil der Frontend-Konzeption. Die Umschaltung erfolgt über einen
dedizierten \passthrough{\lstinline!ModeSelector!} mit Portal-basiertem
Bestätigungsdialog, visueller Toast-Rückmeldung (2200 ms) und
Tastatur-Shortcuts (\passthrough{\lstinline!Ctrl+1!} /
\passthrough{\lstinline!Ctrl+2!} / \passthrough{\lstinline!Ctrl+3!}). Im
kooperativen Modus sind zusätzliche Tastatur-Interaktionen für den
Accept-/Reject-Flow (\passthrough{\lstinline!TAB!} /
\passthrough{\lstinline!ESC!}) eingebunden, um den Redaktionsdurchsatz
zu erhöhen.

\begin{center}\rule{0.5\linewidth}{0.5pt}\end{center}

\subsubsection{4.6.5 Kommunikationsmuster im
Frontend}\label{kommunikationsmuster-im-frontend}

Das Frontend setzt die in Kapitel 4.1.1 definierte hybride
Triggerarchitektur mit drei spezialisierten Mechanismen um:
\textbf{REST-Polling} (5-Sekunden-Intervall) für Kerndaten wie
Ticker-Einträge und Spielevents, \textbf{Webhook-Trigger} für
bedarfsgesteuerte n8n-Importe (vgl. Kap. 4.4.1) sowie \textbf{WebSocket}
(\passthrough{\lstinline!/ws/media!}) mit Exponential-Backoff-Reconnect
für latenzkritische Medieninhalte. REST-Polling wurde gegenüber SSE
gewählt, da die zustandslose HTTP-Architektur die Skalierbarkeit auf
Render vereinfacht. Die konkreten Hook-Implementierungen und
Reconnect-Parameter sind in Kapitel 5.4.3--5.4.5 dokumentiert.

\begin{center}\rule{0.5\linewidth}{0.5pt}\end{center}

\subsection{4.7 Skalierbarkeit, Erweiterbarkeit und
Systemgrenzen}\label{skalierbarkeit-erweiterbarkeit-und-systemgrenzen}

\subsubsection{4.7.1 Horizontale Skalierung des
Backends}\label{horizontale-skalierung-des-backends}

Die Anwendungsschicht ist zustandslos konzipiert und kann horizontal
skaliert werden, indem zusätzliche Uvicorn-Prozesse hinter einem
Load-Balancer gestartet werden. Der gemeinsame
PostgreSQL-Connection-Pool (\passthrough{\lstinline!QueuePool!} mit
\passthrough{\lstinline!pool\_size=20!},
\passthrough{\lstinline!max\_overflow=30!}) stellt sicher, dass mehrere
Backend-Instanzen dieselbe Datenbankverbindungskapazität teilen. Die
WebSocket-Verbindungen für Medien stellen eine Skalierungseinschränkung
dar: Der in-Memory-\passthrough{\lstinline!MediaConnectionManager!}
funktioniert nicht über Prozessgrenzen hinweg und müsste bei
Multi-Prozess-Deployment durch ein verteiltes Pub/Sub-System (z. B.
Redis) ergänzt werden.

\begin{center}\rule{0.5\linewidth}{0.5pt}\end{center}

\subsubsection{4.7.2 Erweiterung um neue
Ereignistypen}\label{erweiterung-um-neue-ereignistypen}

Neue Ereignistypen können ohne Datenbankschema-Änderungen ergänzt
werden: Die \passthrough{\lstinline!events!}-Tabelle verfügt über ein
generisches \passthrough{\lstinline!description!}-Feld, und die
\passthrough{\lstinline!synthetic\_events!}-Tabelle nutzt JSONB für
vollständig flexible Payloads. Die Übersetzung von Football-API-Codes
auf interne Typbezeichnungen ist im Backend in einer zentralen
\passthrough{\lstinline!EVENT\_TYPE\_MAP!}-Konfiguration hinterlegt. Im
Frontend bildet die Funktion \passthrough{\lstinline!getEventMeta()!}
jeden Ereignistyp auf Icon und CSS-Klasse ab --- neue Typen erfordern
ausschließlich einen zusätzlichen Eintrag in dieser Funktion. Alle
übrigen Komponenten konsumieren das normalisierte Icon und sind damit
von konkreten Typ-Strings entkoppelt.

\begin{center}\rule{0.5\linewidth}{0.5pt}\end{center}

\subsubsection{4.7.3 Erweiterung um neue
LLM-Anbieter}\label{erweiterung-um-neue-llm-anbieter}

Der LLM-Dienst dispatcht über ein zentrales Dictionary auf den
providerabhängigen Handler: Für jeden Anbieter existiert eine interne
Methode (\passthrough{\lstinline!\_generate\_openai\_text!},
\passthrough{\lstinline!\_generate\_gemini\_text!} etc.), die im
dispatch-Dictionary unter dem Providernamen registriert ist. Neue
Anbieter können durch Implementierung eines gleichnamigen Handlers und
einen Eintrag in diesem Dictionary ergänzt werden, ohne bestehenden Code
zu berühren. OpenRouter wird bereits als generischer
Proxy-Einstiegspunkt genutzt, der unter Verwendung des
OpenAI-kompatiblen SDK Zugang zu einer Vielzahl weiterer Modelle ohne
separate API-Clients bietet.

\begin{center}\rule{0.5\linewidth}{0.5pt}\end{center}

\subsubsection{4.7.4
Deployment-Architektur}\label{deployment-architektur}

Das System wird auf der Cloud-Plattform \textbf{Render} betrieben und
gliedert sich in drei Deployment-Einheiten:

\begin{enumerate}
\def\labelenumi{\arabic{enumi}.}

\item
  \textbf{Backend} --- Docker-Container
  (\passthrough{\lstinline!python:3.12-slim!}) als Render Web Service.
  Beim Start werden automatisch Datenbankmigrationen ausgeführt, bevor
  der Uvicorn-Server Anfragen entgegennimmt.
\item
  \textbf{Frontend} --- Statisches Build-Artefakt (Create React App),
  das als Render Static Site ausgeliefert wird. Die Konfiguration
  erfolgt zur Buildzeit über Umgebungsvariablen.
\item
  \textbf{PostgreSQL} --- Render Managed Database als zentrale
  Persistenzschicht, auf die sowohl Backend als auch n8n zugreifen.
\end{enumerate}

n8n wird als separater Self-Hosting-Dienst betrieben und kommuniziert
ausschließlich über HTTP-Endpunkte mit dem Backend. Die n8n-Instanz wird
im Entwicklungs- und Evaluationsbetrieb lokal betrieben und via
ngrok-Tunnel für externe Webhooks erreichbar gemacht. Diese
Containerisierung gewährleistet Reproduzierbarkeit zwischen
Entwicklungs- und Produktionsumgebung und ermöglicht eine unabhängige
Skalierung der einzelnen Dienste.

\begin{center}\rule{0.5\linewidth}{0.5pt}\end{center}

\subsection{4.8 Fazit der
Systemkonzeption}\label{fazit-der-systemkonzeption}

Die Systemkonzeption legt fünf interdependente Designpfeiler fest, die
gemeinsam die Produktionsfähigkeit des Systems begründen. Die
dreischichtige Architektur (Kap. 4.1) schafft die strukturelle Basis für
eine unabhängige Weiterentwicklung der Automatisierungs-, Anwendungs-
und Präsentationsschicht. Das relationale Datenbankschema mit seinem
definierten Ticker-Lifecycle (Kap. 4.3) sichert referenzielle Integrität
und Nachvollziehbarkeit aller redaktionellen Entscheidungen --- auch für
spätere Qualitätsanalysen. Die n8n-Workflow-Schicht (Kap. 4.4)
entkoppelt die externe Datenbeschaffung und Orchestrierungslogik vom
Anwendungskern und erlaubt eine schnelle Anpassung von Importprozessen
ohne Eingriff in das Backend. Die Multi-Provider-LLM-Architektur mit
Few-Shot-Prompting und instanzspezifischen Stilprofilen (Kap. 4.5)
maximiert Textqualität und Anbieterunabhängigkeit bei minimaler
Infrastrukturkomplexität. Das Context-basierte Frontend-Design in
Kombination mit dem \passthrough{\lstinline!ticker\_mode!}-Feld (Kap.
4.3.3) ermöglicht redaktionelle Kontrolle ohne Latenzeinbußen und ohne
Prop-Drilling über Komponentengrenzen.

Das vorliegende Konzept ist damit als vollständige Spezifikation für die
Implementierung formuliert, die Kapitel 5 dokumentiert; die
systematische Evaluation folgt in Kapitel 6.
